% ============================================================
%  يُترجَم بـ XeLaTeX  (في Overleaf: Menu → Compiler → XeLaTeX)
%  الخط: Amiri (متوفر في Overleaf، أو حمّله من github.com/aliftype/amiri)
% ============================================================
\documentclass[12pt,a4paper]{article}
\usepackage[a4paper,margin=1.7cm,top=2.3cm,bottom=2.9cm,headheight=15pt,footskip=1.5cm]{geometry}
\usepackage{amsmath,amssymb}
\usepackage{icomma}
\usepackage[english,bidi=default,layout=lists]{babel}
\babelprovide[import,main]{arabic}
\babelfont{rm}[Scale=1.08]{Amiri}
\babelfont{sf}[Scale=1.08]{Amiri}
\usepackage[most]{tcolorbox}
\usepackage{tikz}
\usetikzlibrary{arrows.meta,calc,shadings}
\usepackage{array}
\usepackage{colortbl}
\usepackage{fancyhdr}
\usepackage{titlesec}
\usepackage{enumitem}
\usepackage{collcell}
\usepackage{needspace}
\usepackage{etoolbox}
\newcolumntype{C}{>{\collectcell\ar}c<{\endcollectcell}}
\pretocmd{\section}{\Needspace{15\baselineskip}}{}{}
\pretocmd{\subsection}{\Needspace{11\baselineskip}}{}{}

\definecolor{cDef}{HTML}{1F4E9C}
\definecolor{cThm}{HTML}{B3261E}
\definecolor{cRem}{HTML}{2E7D32}
\definecolor{cEx}{HTML}{C76B00}
\definecolor{cMeth}{HTML}{6A1B9A}
\definecolor{cSol}{HTML}{00695C}
\definecolor{cSec}{HTML}{1F3A68}

\titleformat{\section}{\Large\bfseries\color{cSec}}{\thesection.}{0.6em}{}[{\color{cSec!60}\titlerule[1pt]}]
\titleformat{\subsection}{\large\bfseries\color{cThm!85!black}}{\thesubsection.}{0.5em}{}
\titlespacing*{\section}{0pt}{14pt}{8pt}
\titlespacing*{\subsection}{0pt}{10pt}{5pt}
\setlength{\parindent}{0pt}
\setlength{\parskip}{4pt}
\setlist{itemsep=2pt,topsep=3pt}

\newtcolorbox{arbox}[3][]{enhanced,
  colback=#2!5,colframe=#2!80!black,arc=2mm,boxrule=0.8pt,
  left=6pt,right=6pt,top=8pt,bottom=5pt,
  fonttitle=\bfseries,coltitle=white,
  title={\foreignlanguage{arabic}{#3}},
  attach boxed title to top right={xshift=-5mm,yshift=-3mm},
  boxed title style={colback=#2!80!black,arc=1.5mm,boxrule=0pt},
  before upper={\selectlanguage{arabic}},
  before skip=8pt,after skip=8pt,breakable,#1}
\newenvironment{BOX}[2]{\begin{otherlanguage}{english}\begin{arbox}{#1}{#2}}{\end{arbox}\end{otherlanguage}}
\newenvironment{tarif}[1][تعريف]{\begin{BOX}{cDef}{#1}}{\end{BOX}}
\newenvironment{khassia}[1][خاصية]{\begin{BOX}{cThm}{#1}}{\end{BOX}}
\newenvironment{molahada}[1][ملاحظة]{\begin{BOX}{cRem}{#1}}{\end{BOX}}
\newenvironment{mithal}[1][مثال]{\begin{BOX}{cEx}{#1}}{\end{BOX}}
\newenvironment{tariqa}[1][طريقة]{\begin{BOX}{cMeth}{#1}}{\end{BOX}}

\newcounter{exo}
\newenvironment{tamrin}[1]{\stepcounter{exo}\begin{BOX}{cSec}{التمرين \theexo: #1}}{\end{BOX}}
\newenvironment{hal}{\Needspace{12\baselineskip}\begin{otherlanguage}{english}\begin{arbox}[breakable]{cSol}{حل التمرين \theexo}}{\end{arbox}\end{otherlanguage}}

\newenvironment{LTR}{\begin{otherlanguage}{english}}{\end{otherlanguage}}
\newcommand{\ar}[1]{\foreignlanguage{arabic}{#1}}

\renewcommand{\dots}{\mathinner{\ldotp\ldotp\ldotp}}
\renewcommand{\ldots}{\mathinner{\ldotp\ldotp\ldotp}}
\renewcommand{\cdots}{\mathinner{\cdotp\cdotp\cdotp}}
\newcommand{\N}{\mathbb{N}}
\newcommand{\Z}{\mathbb{Z}}
\newcommand{\D}{\mathbb{D}}
\newcommand{\Q}{\mathbb{Q}}
\newcommand{\R}{\mathbb{R}}
\newcommand{\pgcd}{\operatorname{PGCD}}
\newcommand{\ppcm}{\operatorname{PPCM}}

% \NLseg{a}{b}{[}{]} ; \NLright{a}{[} ; \NLleft{b}{]}  (interval on a number line)
\newcommand{\NLseg}[4]{\begin{LTR}\begin{tikzpicture}[baseline=-0.6ex]
 \draw[-{Stealth[length=2mm]},gray!70](0,0)--(3.4,0);
 \draw[line width=2pt,cThm!80!black](0.7,0)--(2.7,0);
 \node[cThm!80!black,font=\LARGE] at (0.7,0){$#3$};
 \node[cThm!80!black,font=\LARGE] at (2.7,0){$#4$};
 \node[below=3pt,font=\footnotesize] at (0.7,0){$#1$};
 \node[below=3pt,font=\footnotesize] at (2.7,0){$#2$};
\end{tikzpicture}\end{LTR}}
\newcommand{\NLright}[2]{\begin{LTR}\begin{tikzpicture}[baseline=-0.6ex]
 \draw[-{Stealth[length=2mm]},gray!70](0,0)--(3.4,0);
 \draw[line width=2pt,cThm!80!black](1.3,0)--(3.3,0);
 \node[cThm!80!black,font=\LARGE] at (1.3,0){$#2$};
 \node[below=3pt,font=\footnotesize] at (1.3,0){$#1$};
\end{tikzpicture}\end{LTR}}
\newcommand{\NLleft}[2]{\begin{LTR}\begin{tikzpicture}[baseline=-0.6ex]
 \draw[-{Stealth[length=2mm]},gray!70](0,0)--(3.4,0);
 \draw[line width=2pt,cThm!80!black](0.05,0)--(2.1,0);
 \node[cThm!80!black,font=\LARGE] at (2.1,0){$#2$};
 \node[below=3pt,font=\footnotesize] at (2.1,0){$#1$};
\end{tikzpicture}\end{LTR}}

% ---------- social-media icons (pure TikZ) ----------
\definecolor{cFB}{HTML}{1877F2}
\definecolor{cTG}{HTML}{229ED9}
\definecolor{cIG1}{HTML}{FEDA75}
\definecolor{cIG2}{HTML}{D62976}
\definecolor{cIG3}{HTML}{4F5BD5}
\newcommand{\fbicon}{\begin{LTR}\begin{tikzpicture}[baseline=-0.75ex,x=1em,y=1em]
 \fill[cFB] (0,0) circle (0.55);
 \node[white,inner sep=0,font=\usefont{TU}{lmss}{bx}{n}\small] at (0.04,-0.02) {f};
\end{tikzpicture}\end{LTR}}
\newcommand{\igicon}{\begin{LTR}\begin{tikzpicture}[baseline=-0.75ex,x=1em,y=1em]
 \shade[lower left=cIG1,lower right=cIG2,upper left=cIG2,upper right=cIG3,rounded corners=0.28em] (-0.55,-0.55) rectangle (0.55,0.55);
 \draw[white,line width=0.09em,rounded corners=0.18em] (-0.36,-0.36) rectangle (0.36,0.36);
 \draw[white,line width=0.09em] (0,0) circle (0.17);
 \fill[white] (0.22,0.22) circle (0.05);
\end{tikzpicture}\end{LTR}}
\newcommand{\tgicon}{\begin{LTR}\begin{tikzpicture}[baseline=-0.75ex,x=1em,y=1em]
 \fill[cTG] (0,0) circle (0.55);
 \fill[white] (-0.33,0.0) -- (0.32,0.27) -- (0.2,-0.3) -- (0.0,-0.15) -- (-0.08,-0.28) -- (-0.08,-0.08) -- cycle;
 \fill[cTG!70!black] (-0.08,-0.08) -- (0.25,0.2) -- (0.0,-0.15) -- (-0.08,-0.28) -- cycle;
\end{tikzpicture}\end{LTR}}
% page number: black box, white two-digit number, at the right of the social row
\newcommand{\PageBox}{\begin{LTR}\tikz[baseline=(p.base)]\node[fill=black,text=white,inner sep=2.5pt,minimum width=1.6em,minimum height=1.3em,font=\usefont{TU}{lmss}{bx}{n}\small](p){\ifnum\value{page}<10 0\fi\arabic{page}};\end{LTR}}
% visual left-to-right order: Telegram | Instagram | Facebook | page box
\newcommand{\SocialRow}{\begin{otherlanguage}{english}%
 \ar{\textbf{الأستاذ سوايسية هشام للرياضيات}}\ \tgicon\hfill
 \textbf{prof\_hicham\_math}\ \igicon\hfill
 \ar{\textbf{الأستاذ سوايسية هشام للرياضيات}}\ \fbicon\hspace{2.2em}\PageBox\end{otherlanguage}}

\pagestyle{fancy}
\fancyhf{}
\renewcommand{\headrulewidth}{0.6pt}
\renewcommand{\footrulewidth}{0.4pt}
% \SetHeader{right head: document title}{left head: level -- stream}
\newcommand{\SetHeader}[2]{%
 \fancyhead[C]{\parbox{\textwidth}{\ar{#1\hfill #2}}}%
 \fancyfoot[C]{\parbox{\textwidth}{\small\SocialRow}}}
\usepackage{tkz-tab}
\usepackage{multicol}
% variation / sign table (LTR)
\newenvironment{TAB}{\par\begin{center}\begin{LTR}\begin{tikzpicture}[scale=0.95,every node/.style={font=\small}]}{\end{tikzpicture}\end{LTR}\end{center}}
% grid + axes, clipped: \begin{GR}{xmin}{xmax}{ymin}{ymax}{unit cm}
\newenvironment{GR}[5]{\par\begin{center}\begin{LTR}\begin{tikzpicture}[x=#5cm,y=#5cm,font=\footnotesize]
 \draw[gray!25,very thin] (#1,#3) grid (#2,#4);
 \draw[-{Stealth[length=2mm]}] (#1,0)--(#2+0.3,0) node[below left]{$x$};
 \draw[-{Stealth[length=2mm]}] (0,#3)--(0,#4+0.3) node[below left]{$y$};
 \node[below left] at (0,0){$O$};
 \begin{scope}\clip (#1,#3) rectangle (#2,#4);}
{\end{scope}\end{tikzpicture}\end{LTR}\end{center}}
\newcommand{\Ticks}[2]{\foreach \t in {#1} \node[below,font=\scriptsize] at (\t,0){$\t$}; \foreach \t in {#2} \node[left,font=\scriptsize] at (0,\t){$\t$};}
\newcommand{\bx}[1]{\ensuremath{\colorbox{cDef!12}{$\displaystyle #1$}}}
\newcommand{\TwoCol}[4]{\par\begin{LTR}\noindent\begin{minipage}[t]{#1\textwidth}\selectlanguage{arabic}#2\end{minipage}\hfill\begin{minipage}[t]{#3\textwidth}\selectlanguage{arabic}#4\end{minipage}\end{LTR}\par}
\newcommand{\RL}[4]{\TwoCol{#3}{#4}{#1}{#2}}  % \RL{w1}{right content}{w2}{left content}
\newcommand{\ok}{$\textcolor{cRem}{\checkmark}$}
\newcommand{\limn}{\displaystyle\lim_{n\to+\infty}}
% sub-questions, each on its own line
\newenvironment{sq}{\begin{itemize}[leftmargin=2.4em,itemsep=3pt]}{\end{itemize}}
\newcommand{\qa}{\item[(أ)]}\newcommand{\qb}{\item[(ب)]}\newcommand{\qc}{\item[(جـ)]}\newcommand{\qd}{\item[(د)]}\newcommand{\qe}{\item[(هـ)]}
\newcommand{\pa}{\par\smallskip\textbf{(أ)} }\newcommand{\pb}{\par\smallskip\textbf{(ب)} }\newcommand{\pc}{\par\smallskip\textbf{(جـ)} }\newcommand{\pd}{\par\smallskip\textbf{(د)} }\newcommand{\pe}{\par\smallskip\textbf{(هـ)} }
\newcommand{\QQ}[1]{\par\smallskip\textbf{#1.} }
\newenvironment{tmr}{\stepcounter{exo}\begin{BOX}{cSec}{تمرين محلول \theexo}}{\end{BOX}}
\newcommand{\ABC}[3]{\\[2pt]\begin{tabular}{p{0.28\textwidth}p{0.28\textwidth}p{0.28\textwidth}}(جـ) \ #3 & (ب) \ #2 & (أ) \ #1\end{tabular}}
\pretocmd{\section}{\Needspace{24\baselineskip}}{}{}
\newtcbox{\bxtbox}{on line,colback=cDef!12,colframe=cDef!12,boxrule=0pt,arc=1pt,boxsep=0pt,left=2pt,right=2pt,top=1pt,bottom=1pt}
\newcommand{\bxt}[1]{\foreignlanguage{english}{\bxtbox{\ar{#1}}}}
\usepackage{longtable}
\setlength{\LTpre}{4pt}\setlength{\LTpost}{0pt}
\newcommand{\Phase}[1]{\textbf{#1}}
\newcommand{\Fld}[2]{\textbf{#1:}\ #2}
% رأس المذكرة على نمط النموذج: إطار بعنوان «مذكرة رقم»، معلومات الحصة، ثم إطار الكفاءات
\newcommand{\ModhHead}[9]{%
\begin{LTR}
\begin{tcolorbox}[enhanced,colback=white,colframe=cThm!85!black,arc=4mm,boxrule=1pt,top=14pt,bottom=6pt,left=8pt,right=8pt,
 overlay={\node[fill=cThm!85!black,text=white,rounded corners=3mm,inner xsep=14pt,inner ysep=4pt,font=\large\bfseries] at ([yshift=0pt]frame.north){\ar{مذكرة رقم #1}};},
 before upper={\selectlanguage{arabic}}]
\begin{tabular}{@{}p{0.4\linewidth}p{0.55\linewidth}@{}}
\Fld{المستوى}{#2} & \Fld{الثانوية}{\dotfill}\\[3pt]
\Fld{المادة}{رياضيات} & \Fld{ميدان التعلم}{#3}\\[3pt]
\Fld{المدة}{#5} & \Fld{الوحدة التعلمية}{#4}\\[3pt]
\Fld{الأستاذ}{\dotfill} & \Fld{المحتوى المعرفي}{#6}\\
\end{tabular}
\end{tcolorbox}
\begin{tcolorbox}[enhanced,colback=white,colframe=cThm!85!black,arc=3mm,boxrule=0.8pt,sidebyside,sidebyside align=top,lefthand ratio=0.5,
 segmentation style={dashed,cThm!60},before upper={\selectlanguage{arabic}},before lower={\selectlanguage{arabic}}]
\textbf{الكفاءات المستهدفة:}\par #8
\tcblower
\textbf{الكفاءات القبلية:}\par #7
\end{tcolorbox}
\end{LTR}
\par\smallskip\textbf{الوسائل:} #9\par\smallskip}
\newenvironment{Sayr}{%
\begin{LTR}\renewcommand{\arraystretch}{1.25}
\begin{longtable}{|p{0.17\textwidth}|>{\centering\arraybackslash}p{0.06\textwidth}|p{0.53\textwidth}|>{\centering\arraybackslash}p{0.12\textwidth}|}
\hline\rowcolor{cThm!12}\ar{\textbf{توجيهات وتعاليق}} & \ar{\textbf{المدة}} & \ar{\textbf{سير الحصة}} & \ar{\textbf{مراحل الحصة}}\\\hline\endhead
\hline\endfoot}{\end{longtable}\end{LTR}}
% سطر في الجدول: \Row{مرحلة}{سير}{مدة}{توجيهات}
\newcommand{\Row}[4]{\ar{\small #4} & \ar{#3} & \ar{#2} & \cellcolor{cThm!8}\ar{\Phase{#1}}\\\hline}

\newcommand{\Lab}[2]{\foreignlanguage{english}{\tcbox[on line,colback=#1,colframe=#1,boxrule=0pt,arc=2pt,boxsep=1pt,left=5pt,right=5pt,top=1pt,bottom=1pt]{\ar{\textcolor{white}{\textbf{#2}}}}}}
\newcommand{\Num}[1]{\foreignlanguage{english}{\tikz[baseline=(c.base)]\node[fill=cDef,text=white,inner sep=2pt,font=\footnotesize\bfseries](c){#1};}}
\newcommand{\ck}{\ensuremath{{\color{cRem}\checkmark}}}
\newenvironment{Kad}[1]{\par\smallskip\Lab{#1}}{\par}
\SetHeader{\small\bfseries\color{cSec}مذكرات محور: النهايات}{\small\color{cSec}السنة الثالثة ثانوي -- علوم تجريبية}
\begin{document}
\begin{LTR}
\begin{tcolorbox}[enhanced,colback=cThm!85!black,colframe=cThm!85!black,arc=3mm,halign=center,boxrule=0pt,top=12pt,bottom=12pt,before upper={\selectlanguage{arabic}}]
{\color{white}\fontsize{26}{32}\selectfont\bfseries مذكرات محور: النهايات}\\[8pt]
{\color{white}\large السنة الثالثة ثانوي -- علوم تجريبية (صالحة لشعبتي رياضيات وتقني رياضي)}\\[4pt]
{\color{white}\large 5 مذكرات (10 ساعات) -- السنة الدراسية 2026-2027}
\end{tcolorbox}
\end{LTR}
\begin{BOX}{cThm}{يتضمن هذا العمل}
\begin{enumerate}[label=\textbf{مذكرة \arabic{enumi}:},leftmargin=6em,start=18]
\item النهاية عند $\pm\infty$ ، المستقيم المقارب الموازي لمحور الفواصل.
\item النهاية عند عدد حقيقي ، المستقيم المقارب الموازي لمحور التراتيب.
\item العمليات على النهايات وإزالة حالات عدم التعيين.
\item النهايات بالمقارنة والحصر ، نهاية دالة مركبة.
\item السلوك التقاربي: المستقيم المقارب المائل.
\end{enumerate}
مراجع التمارين: الكتاب المدرسي للسنة الثالثة ثانوي ، الجزء الأول (الباب الأول: النهايات والاستمرارية ص 26--27). التزايد المقارن يأتي في المحور الموالي.
\end{BOX}
\newpage
\ModhHead{18}{3 ع ت}{تحليل}{النهايات}{02 سا}{النهاية عند $\pm\infty$ ، المستقيم المقارب الموازي لمحور الفواصل}%
{\begin{itemize}\item النهايات (2 ع ت) ، الدالتان الأسية واللوغاريتمية.\end{itemize}}%
{\begin{itemize}\item حساب نهاية دالة عند $\pm\infty$ ، والتفسير البياني.\end{itemize}}%
{المنهاج ، الكتاب المدرسي ، مراجع ، الحاسبة البيانية.}
\begin{Sayr}
\Row{مرحلة الانطلاق}{%
\textbf{التهيئة النفسية:} ذكّر بنهايات $e^x$ و $\ln x$ عند حدود مجموعتي تعريفهما.}{10 د}{تذكير.}
\Row{التشخيص والاكتشاف}{%
\Lab{cEx}{وضعية مقترحة}\ $f(x)=2+e^{-x}$ و $g(x)=\dfrac{3x-1}{x+2}$: جداول قيم لما $x$ كبير جداً.\par\ck\ القيم تقترب من $2$ ومن $3$: المنحنيان يقتربان من مستقيمين أفقيين.\par
}{35 د}{جداول قيم.}
\Row{بناء المعارف}{%
\Lab{cDef}{تعريف}\ إذا كان $\lim\limits_{x\to +\infty}f(x)=l$ (أو عند $-\infty$) فالمستقيم $y=l$ \textbf{مقارب} لـ $C_f$ بجوار $+\infty$.\par\Lab{cThm}{نهايات مرجعية}\ $\lim\limits_{x\to +\infty}x^n=+\infty$ ؛ $\lim\limits_{x\to \pm\infty}\dfrac1{x^n}=0$ ؛ $\lim\limits_{x\to +\infty}e^x=+\infty$ ، $\lim\limits_{x\to -\infty}e^x=0$ ؛ $\lim\limits_{x\to +\infty}\ln x=+\infty$.\par}{50 د}{جدول النهايات المرجعية.}
\Row{}{%
\Lab{cEx}{مثال}\ $\lim\limits_{x\to +\infty}\left(2+e^{-x}\right)=2$: $y=2$ مقارب ؛ $\lim\limits_{x\to \pm\infty}\dfrac{3x-1}{x+2}=3$: $y=3$ مقارب.\par\Lab{cThm!80!black}{حذار}\ نحدد الجهة: المقارب قد يكون عند $+\infty$ فقط.}{}{}
\Row{التقويم}{%
\Lab{cEx}{تطبيق}\par احسب: $\lim\limits_{x\to -\infty}(1+e^x)$ ؛ $\lim\limits_{x\to +\infty}\dfrac1{\ln x}$ ، وفسّر بيانياً.\par\textbf{الحل:} $1$ ($y=1$ مقارب بجوار $-\infty$) ؛ $0$ ($y=0$ مقارب بجوار $+\infty$).\par\textbf{واجب منزلي:} تمارين من الباب الأول.\par\smallskip\textbf{ملاحظات حول سير الحصة:} \dotfill}{25 د}{تقويم فردي.}
\end{Sayr}
\newpage
\ModhHead{19}{3 ع ت}{تحليل}{النهايات}{02 سا}{النهاية عند عدد حقيقي ، المستقيم المقارب الموازي لمحور التراتيب}%
{\begin{itemize}\item النهايات عند $\pm\infty$.\end{itemize}}%
{\begin{itemize}\item حساب نهاية دالة عند حد منته لمجموعة تعريفها ، والتفسير البياني.\end{itemize}}%
{المنهاج ، الكتاب المدرسي ، مراجع ، الحاسبة البيانية.}
\begin{Sayr}
\Row{مرحلة الانطلاق}{%
\textbf{التهيئة النفسية:} ماذا يحدث لـ $\dfrac1{x-2}$ لما يقترب $x$ من $2$ من اليمين؟ من اليسار؟}{10 د}{سؤال.}
\Row{التشخيص والاكتشاف}{%
\Lab{cEx}{وضعية مقترحة}\par\ck\ جداول قيم لـ $\dfrac{x+1}{x-2}$ من اليمين ومن اليسار.\par
\ck\ نلاحظ نهايتين غير منتهيتين مختلفتي الإشارة.\par
}{35 د}{جداول اليمين واليسار.}
\Row{بناء المعارف}{%
\Lab{cThm}{خاصية}\ إذا كانت $\lim\limits_{x\to a}f(x)=\pm\infty$ (من جهة على الأقل) فالمستقيم $x=a$ مقارب لـ $C_f$.\par\Lab{cMeth}{طريقة}\ البسط يؤول إلى عدد غير معدوم والمقام إلى $0$: ندرس إشارة المقام ($0^+$ أو $0^-$).\par}{50 د}{جدول إشارة المقام.}
\Row{}{%
\Lab{cEx}{مثال}\par\ck\ $\lim\limits_{x\to 2^+}\dfrac{x+1}{x-2}=+\infty$ ، $\lim\limits_{x\to 2^-}\dfrac{x+1}{x-2}=-\infty$: $x=2$ مقارب.\par
\ck\ $\lim\limits_{x\to 1^+}\dfrac1{\ln x}=+\infty$ ، $\lim\limits_{x\to 1^-}\dfrac1{\ln x}=-\infty$: $x=1$ مقارب.\par
\ck\ $\lim\limits_{x\to 0^+}\ln x=-\infty$: محور التراتيب مقارب لمنحني $\ln$.\par
\Lab{cThm!80!black}{حذار}\ $\dfrac{k}{0}$ ليست نتيجة: نحدد الإشارة.}{}{}
\Row{التقويم}{%
\Lab{cEx}{تطبيق}\par احسب نهاية $h(x)=\dfrac{x}{(x+3)^2}$ عند $-3$ وفسّر بيانياً.\par\textbf{الحل:} البسط $\to-3$ والمقام $\to0^+$: $-\infty$ من الجهتين ؛ $x=-3$ مقارب.\par\textbf{واجب منزلي:} تمارين 13 ، 14 ، 16 صفحة 27.\par\smallskip\textbf{ملاحظات حول سير الحصة:} \dotfill}{25 د}{تقويم فردي.}
\end{Sayr}
\newpage
\ModhHead{20}{3 ع ت}{تحليل}{النهايات}{02 سا}{العمليات على النهايات وإزالة حالات عدم التعيين}%
{\begin{itemize}\item النهايات المرجعية ، العدد المشتق.\end{itemize}}%
{\begin{itemize}\item حساب نهايات بمبرهنات العمليات ، وإزالة حالات عدم التعيين.\end{itemize}}%
{المنهاج ، الكتاب المدرسي ، مراجع ، الحاسبة البيانية.}
\begin{Sayr}
\Row{مرحلة الانطلاق}{%
\textbf{التهيئة النفسية:} ما حالات عدم التعيين الأربع؟}{10 د}{تذكير.}
\Row{التشخيص والاكتشاف}{%
\Lab{cEx}{وضعية مقترحة}\par\ck\ $\lim\limits_{x\to 0}\dfrac{e^x-1}x$ من الشكل $\frac00$: نتعرف فيه على نسبة تزايد الدالة $\exp$ عند $0$.\par
}{35 د}{التعرف على عدد مشتق.}
\Row{بناء المعارف}{%
\Lab{cThm}{تذكير}\ مبرهنات المجموع والجداء والقسمة ، وحالات عدم التعيين: $\infty-\infty$ ، $0\times\infty$ ، $\dfrac\infty\infty$ ، $\dfrac00$.\par\Lab{cMeth}{إزالة عدم التعيين}\par\ck\ كثير حدود وناطقة عند $\pm\infty$: الحد الأعلى درجة.\par
\ck\ جذور: الضرب في المرافق.\par
\ck\ $\frac00$ عند $a$: التحليل ، أو التعرف على عدد مشتق: $\lim\limits_{x\to a}\dfrac{f(x)-f(a)}{x-a}=f^{\prime}(a)$.\par
}{50 د}{جدول الطرق.}
\Row{}{%
\Lab{cEx}{مثال}\par\ck\ $\lim\limits_{x\to +\infty}\left(\sqrt{x^2+1}-x\right)=\lim\limits_{x\to +\infty}\dfrac1{\sqrt{x^2+1}+x}=0$.\par
\ck\ $\lim\limits_{x\to 0}\dfrac{e^x-1}x=\exp^{\prime}(0)=1$ ؛ $\lim\limits_{x\to 0}\dfrac{\ln(1+x)}x=1$.\par
\Lab{cThm!80!black}{حذار}\ «$\infty-\infty$» لا تساوي $0$.}{}{}
\Row{التقويم}{%
\Lab{cEx}{تطبيق}\par احسب: $\lim\limits_{x\to +\infty}\dfrac{2x^2-x}{x^2+3}$ ؛ $\lim\limits_{x\to 0}\dfrac{\ln(1+2x)}x$.\par\textbf{الحل:} $2$ ؛ العدد المشتق لـ $\ln(1+2x)$ عند $0$: $2$.\par\textbf{واجب منزلي:} تمارين 18 ، 19 ، 24 صفحة 26.\par\smallskip\textbf{ملاحظات حول سير الحصة:} \dotfill}{25 د}{تقويم فردي.}
\end{Sayr}
\newpage
\ModhHead{21}{3 ع ت}{تحليل}{النهايات}{02 سا}{النهايات بالمقارنة والحصر ، نهاية دالة مركبة}%
{\begin{itemize}\item العمليات على النهايات.\end{itemize}}%
{\begin{itemize}\item حساب نهايات باستعمال المقارنة أو الحصر أو تركيب دالتين.\end{itemize}}%
{المنهاج ، الكتاب المدرسي ، مراجع ، الحاسبة البيانية.}
\begin{Sayr}
\Row{مرحلة الانطلاق}{%
\textbf{التهيئة النفسية:} $\sin x$ ليست لها نهاية عند $+\infty$: فهل لـ $\dfrac{\sin x}x$ نهاية؟}{10 د}{سؤال مفتوح.}
\Row{التشخيص والاكتشاف}{%
\Lab{cEx}{وضعية مقترحة}\par\ck\ $-\dfrac1x\leqslant\dfrac{\sin x}x\leqslant\dfrac1x$ لما $x>0$ ، والطرفان يؤولان إلى $0$.\par
}{35 د}{رسم بالحاسبة.}
\Row{بناء المعارف}{%
\Lab{cThm}{المقارنة}\ إذا كان $f(x)\geqslant g(x)$ بجوار $+\infty$ و $\lim\limits_{x\to +\infty}g(x)=+\infty$ فإن $\lim\limits_{x\to +\infty}f(x)=+\infty$.\par\Lab{cThm}{الحصر}\ إذا كان $g(x)\leqslant f(x)\leqslant h(x)$ و $g$ ، $h$ لهما نفس النهاية $l$ فإن $\lim\limits_{x\to +\infty}f(x)=l$.\par}{50 د}{الطرق الثلاث.}
\Row{}{%
\Lab{cThm}{دالة مركبة}\ إذا كان $\lim\limits_{x\to a}u(x)=b$ و $\lim\limits_{x\to b}g(X)=c$ فإن $\lim\limits_{x\to a}g\big(u(x)\big)=c$.\par\Lab{cEx}{مثال}\ $x+\sin x\geqslant x-1$: النهاية $+\infty$ ؛ $\lim\limits_{x\to +\infty}e^{-x^2}=0$ لأن $-x^2\to-\infty$ ؛ $\lim\limits_{x\to +\infty}\sqrt{4+\frac1x}=2$.\par}{}{}
\Row{}{%
\Lab{cThm!80!black}{حذار}\ نكتب الحصر بوضوح ونتحقق من نهايتي الطرفين.}{}{}
\Row{التقويم}{%
\Lab{cEx}{تطبيق}\par احسب: $\lim\limits_{x\to +\infty}\dfrac{2+\cos x}x$ ؛ $\lim\limits_{x\to +\infty}e^{\frac1x}$.\par\textbf{الحل:} $\dfrac1x\leqslant\dfrac{2+\cos x}x\leqslant\dfrac3x$: النهاية $0$ ؛ $\frac1x\to0$: النهاية $e^0=1$.\par\textbf{واجب منزلي:} تمارين من الباب الأول.\par\smallskip\textbf{ملاحظات حول سير الحصة:} \dotfill}{25 د}{تقويم فردي.}
\end{Sayr}
\newpage
\ModhHead{22}{3 ع ت}{تحليل}{النهايات}{02 سا}{السلوك التقاربي: المستقيم المقارب المائل}%
{\begin{itemize}\item النهايات ، القسمة الإقليدية.\end{itemize}}%
{\begin{itemize}\item تبرير أن مستقيماً معلوماً مقارب مائل ، والبحث عنه ، ودراسة الوضع النسبي.\end{itemize}}%
{المنهاج ، الكتاب المدرسي ، مراجع ، الحاسبة البيانية.}
\begin{Sayr}
\Row{مرحلة الانطلاق}{%
\textbf{التهيئة النفسية:} منحني $f(x)=x-1+e^{-x}$ بالحاسبة: إلى أي مستقيم يقترب لما $x\to+\infty$ ؟}{10 د}{رسم بالحاسبة.}
\Row{التشخيص والاكتشاف}{%
\Lab{cEx}{وضعية مقترحة}\par\ck\ $f(x)-(x-1)=e^{-x}\to0$ عند $+\infty$ ، وهو موجب: $C_f$ فوق المستقيم.\par
}{35 د}{القسمة الإقليدية.}
\Row{بناء المعارف}{%
\Lab{cDef}{تعريف}\ $y=ax+b$ مقارب مائل لـ $C_f$ بجوار $+\infty$ إذا كان $\lim\limits_{x\to +\infty}\big[f(x)-(ax+b)\big]=0$.\par\Lab{cMeth}{البحث عن المقارب}\ نكتب $f(x)=ax+b+g(x)$ مع $\lim\limits_{x\to \pm\infty}g(x)=0$ (قسمة إقليدية لدالة ناطقة).\par}{50 د}{جدول الإشارة.}
\Row{}{%
\Lab{cEx}{مثال}\ $f(x)=\dfrac{x^2+x+1}{x-1}=x+2+\dfrac3{x-1}$: $y=x+2$ مقارب مائل ؛ $C_f$ فوقه لما $x>1$ وتحته لما $x<1$.\par\Lab{cMeth}{الوضع النسبي}\ إشارة $f(x)-(ax+b)$.\par}{}{}
\Row{}{%
\Lab{cThm!80!black}{حذار}\ نهاية الفرق يجب أن تكون $0$ بالضبط.}{}{}
\Row{التقويم}{%
\Lab{cEx}{تطبيق}\par $f(x)=2x+3+\dfrac1{x^2+1}$. بيّن أن $y=2x+3$ مقارب مائل ، وادرس الوضع النسبي.\par\textbf{الحل:} $f(x)-(2x+3)=\dfrac1{x^2+1}\to0$ عند $\pm\infty$ ، وهو موجب: $C_f$ فوق المقارب.\par\textbf{واجب منزلي:} تمارين من الباب الأول (المقارب المائل).\par\smallskip\textbf{ملاحظات حول سير الحصة:} \dotfill}{25 د}{نهاية المحور: فرض أو مراجعة.}
\end{Sayr}
\end{document}
